% !TeX root = ../Book1.tex
% Möbius线丛的端点粘合；纤维只截取有限区间用于示意。
\begin{tikzpicture}[x=0.95cm,y=1.15cm,>=stealth]
	\fill[blue!3] (0,-1.35) rectangle (6.283,1.35);
	\draw[gray!65,dashed] (0,1.35)--(6.283,1.35) (0,-1.35)--(6.283,-1.35);
	\draw[very thick,->] (0,-1.35)--(0,1.35);
	\draw[very thick,->] (6.283,1.35)--(6.283,-1.35);
	\draw[gray!65,->] (-0.3,0)--(6.8,0) node[right] {\(t\)};
	\node[above] at (0,1.4) {\(v\)};
	\node[below left] at (0,0) {\(0\)};
	\node[below] at (3.142,-0.04) {\(\pi\)};
	\node[below right] at (6.283,0) {\(2\pi\)};
	\draw[red!75!black,dashed,thick] (0,0.55)--(6.283,0.55);
	\draw[red!75!black,fill=white] (0,0.55) circle (2pt) (6.283,0.55) circle (2pt);
	\node[red!75!black,above] at (4.9,0.6) {恒正分量不相合};
	\draw[blue!75!black,very thick,domain=0:6.283,samples=80] plot (\x,{cos(deg(\x/2))});
	\fill[blue!75!black] (0,1) circle (2pt) (3.142,0) circle (2pt) (6.283,-1) circle (2pt);
	\node[blue!75!black,above right] at (0.5,1.02) {\(f(t)=\cos(t/2)\)};
	\node[left] at (0,1) {\(1\)};
	\node[right] at (6.283,-1) {\(-1\)};
	\node[below,align=center] at (3.142,-1.52) {\((2\pi,v)\sim(0,-v)\)：两侧沿相反方向粘合};
\end{tikzpicture}
